\subsection{The projective arithmetic source}

Paper I separates literal histories from their induced projective operators.  That
distinction is decisive here.  Over
\[
 K_4=\Q(\sqrt2)
\]
the bilateral projective history language contains the two operators
\begin{equation}
 \mathsf T(\tau)=\tau+\sqrt2,
 \qquad
 \mathsf J(\tau)=-\frac1\tau,
 \label{eq:q4-projective-generators}
\end{equation}
represented by
\[
 T=\begin{pmatrix}1&\sqrt2\\0&1\end{pmatrix},
 \qquad
 J=\begin{pmatrix}0&-1\\1&0\end{pmatrix}.
\]
Projective evaluation therefore contains the operator subgroup
\begin{equation}
 \Gamma_4=\langle J,T\rangle<\operatorname{PSL}_2(\R),
 \qquad J^2=(JT)^4=1.
 \label{eq:q4-hecke-group}
\end{equation}
It is the Hecke triangle group of signature $(2,4,\infty)$.  Its parabolic
translation is $T$, and the elliptic generators have orders $2$ and $4$.
The associated Hecke--Farey and continued-fraction paths are arithmetic paths in
the operator quotient, not distinct merely because they have different word
spellings \cite{Rosen1954ContinuedFractions,LangLang2016HeckeGroups}.

Let
\[
 \beta:\HH\longrightarrow\mathbb P^1(\C)
\]
be the normalized $\Gamma_4$-invariant Hauptmodul.  We choose the normalization
for which the order-$2$ elliptic orbit maps to $0$, the order-$4$ elliptic orbit
maps to $1$, and the cusp orbit maps to $\infty$.  Thus
\begin{equation}
 \Gamma_4\backslash\HH^*\cong\mathbb P^1,
 \qquad
 (\mathcal E_2,\mathcal E_4,\mathcal C)
 \longmapsto(0,1,\infty),
 \label{eq:q4-hauptmodul-normalization}
\end{equation}
with finite ramification indices $2$ and $4$ at the elliptic orbits and a cusp over
$\infty$.  The existence and normalization are standard for the triangle orbifold
\cite{Beardon1983DiscreteGroups,DoranEtAl2013TriangleForms}.  If a cited convention
interchanges the two finite branch values, we replace its Hauptmodul by $1-\beta$;
this is the only renormalization being made.  In particular, in
suitable local coordinates,
\begin{align}
 \beta(z)&=c_2z^2+O(z^3) &&\text{at }\mathcal E_2,\notag\\
 1-\beta(z)&=c_4z^4+O(z^5) &&\text{at }\mathcal E_4,
 \label{eq:q4-elliptic-expansions}
\end{align}
where $c_2c_4\ne0$.  In a quotient cusp coordinate $q$, the Hauptmodul has a
simple pole,
\begin{equation}
 \beta(q)=c_\infty q^{-1}+O(1),
 \qquad c_\infty\ne0.
 \label{eq:q4-cusp-expansion}
\end{equation}

\subsection{The Hecke--Hauptmodul singular AES}

The interval $[0,1]$ is the edge between the two finite orbifold branch values.
The next construction turns its full preimage into an exact AEG zero network.

\begin{theorem}[$q=4$ Hecke zero network]
\label{thm:q4-hecke-zero-network}
Let $\beta$ have the normalization
\eqref{eq:q4-hauptmodul-normalization}, fix $\mu\ne0$ and $\lambda\in\R$, and
choose a meromorphic square root on $\HH$
\begin{equation}
 F^2=\frac{\beta}{1-\beta}.
 \label{eq:q4-square-root}
\end{equation}
Define
\begin{equation}
 W=\frac{F-i}{F+i},
 \qquad
 s=\log|W|,
 \label{eq:q4-cayley-coordinate}
\end{equation}
and
\begin{equation}
 a_4=
 \begin{cases}
 \dfrac{\mu}{\lambda}\sinh(\lambda s),&\lambda\ne0,\\[6pt]
 \mu s,&\lambda=0,
 \end{cases}
 \qquad
 g_4=\left|\frac{\dd W}{W}\right|^2
 \quad\text{on }\HH\setminus\mathcal E_4.
 \label{eq:q4-aeg-data}
\end{equation}
Then:
\begin{enumerate}
\item $(\HH,\mathcal E_4,g_4,a_4;\mu,\lambda)$ is a singular AES, and every
point of $\mathcal E_4$ is locally essential;
\item its total zero set is exactly
\begin{equation}
 Z(a_4)=\beta^{-1}([0,1]);
 \label{eq:q4-zero-dessin}
\end{equation}
\item points of $\mathcal E_2$ are bivalent regular points of the zero set, while
points of $\mathcal E_4$ are four-valent singular zeros; the metric completion at
each latter point has cone angle $4\pi$;
\item after suppressing the bivalent $\mathcal E_2$ vertices,
$\beta^{-1}([0,1])$ is the $\{\infty,4\}$ Hecke zero network: every vertex has
valence four and every face runs out to a parabolic cusp and has infinitely many
sides;
\item the punctured regular metric on $\HH\setminus\mathcal E_4$ is incomplete at
$\mathcal E_4$, whose points are at finite distance.  Its upstairs completion adds
$4\pi$ cone points.  The induced $\Gamma_4$-orbifold length metric, after those
points are adjoined, is complete and has its cusp at infinite distance.  The
underlying coarse quotient has angle $\pi$ at the image of an order-$4$ point;
none of these statements is smooth geodesic completeness across a regular AES
point.
\end{enumerate}
\end{theorem}

\begin{proof}
The divisor of $\beta/(1-\beta)$ on $\HH$ has zeros of order $2$ along
$\mathcal E_2$ and poles of order $4$ along $\mathcal E_4$.  It is therefore
even.  Since $\HH$ is simply connected, it admits the meromorphic square root
\eqref{eq:q4-square-root}, unique up to sign.  The equation $F=\pm i$ would imply
$F^2=-1=\beta/(1-\beta)$, which is impossible.  Hence the Cayley transform $W$
extends holomorphically through the poles of $F$ and is a nowhere-zero finite
holomorphic function on $\HH$.

Away from $\mathcal E_4$, $\dd W$ is nonzero: $\beta$ has no ramification away
from the two elliptic orbits, the square root cancels the order-$2$ ramification
at $\mathcal E_2$, and the Cayley map has nonzero derivative.  Locally choose a
logarithm $\zeta=\log W$.  Equations~\eqref{eq:q4-cayley-coordinate} and
\eqref{eq:q4-aeg-data} become
\[
 s=\re\zeta,
 \qquad g_4=|\dd\zeta|^2.
\]
Thus $|\nabla s|_{g_4}=1$.  If $\lambda\ne0$, then
\[
 \frac{\dd a_4}{\dd s}=\mu\cosh(\lambda s),
 \qquad
 |\nabla a_4|_{g_4}^2
 =\mu^2\cosh^2(\lambda s)
 =\mu^2+\lambda^2a_4^2.
\]
The calculation for $\lambda=0$ is immediate.

The function of $s$ defining $a_4$ is odd and vanishes only at $s=0$.  Moreover,
the Cayley transform carries $\mathbb P^1(\R)$ to the unit circle, so
\begin{align*}
 a_4=0
 &\iff |W|=1
 \iff F\in\mathbb P^1(\R)\\
 &\iff \frac{\beta}{1-\beta}\in[0,\infty]
 \iff \beta\in[0,1].
\end{align*}
This proves \eqref{eq:q4-zero-dessin}.

At a point of $\mathcal E_2$, expansion
\eqref{eq:q4-elliptic-expansions} gives $F=c z+O(z^2)$ and
$W=-1+c'z+O(z^2)$ with $cc'\ne0$.  The zero set is therefore a smooth arc.
At a point of $\mathcal E_4$, one has
$F=c z^{-2}(1+O(z))$ and hence
\[
 W=1+c'z^2+O(z^3),
 \qquad c'\ne0.
\]
Thus $\log W$ has local degree $2$.  Its real zero has four prongs, and the local
metric is the pullback of $|\dd\zeta|^2$ by a degree-$2$ branch.  The cone-angle
calculation in Theorem~\ref{thm:branched-pullback-zero-germ} gives $4\pi$.
The assignment is critical there, so local essentiality follows from the same
$|\nabla a|^2\ge\mu^2$ obstruction.

The graph $\beta^{-1}([0,1])$ is the orbit of the finite edge in a
$(2,4,\infty)$ triangle.  Four such edge germs meet at an order-$4$ vertex; an
order-$2$ point merely subdivides an edge.  Suppressing those subdivisions gives
the standard $\{\infty,4\}$ Hecke graph.  Its complementary regions are the
parabolic cusp regions and have infinitely many boundary edges
\cite{LangLang2016HeckeGroups}.

Finally, Proposition~\ref{prop:q4-sign-descent} below shows that $g_4$ descends
and that $\chi(T)=+1$ for the cusp stabilizer.  Thus $F$ and $W$ are single-valued
in the quotient cusp coordinate $q$.  The cone calculation places every point of $\mathcal E_4$ at finite
distance upstairs, with angle $4\pi$.  Dividing by its order-$4$ stabilizer gives
angle $\pi$ on the underlying coarse quotient.  At a cusp,
\eqref{eq:q4-cusp-expansion} gives
\[
 W=Cq^{\pm1}(1+O(q)),
 \qquad
 \frac{\dd W}{W}=\pm\frac{\dd q}{q}+O(\dd q).
\]
The radial integral $\int|\dd q|/|q|$ diverges.  The completed orbifold quotient
has a compact core after cusp neighborhoods are removed, so adjoining the
finite-distance singular orbit and retaining the infinite cusp end gives the
asserted length-metric completeness.  The order-$2$ orbit is smooth upstairs but
also has coarse quotient angle $\pi$ because of its isotropy; its bivalent zero
germ distinguishes it from the order-$4$ singular orbit.
\end{proof}

\subsection{Sign-character descent}

The square root in \eqref{eq:q4-square-root} is not $\Gamma_4$-invariant.  This
is not an ambiguity to be hidden; it records a natural signed assignment.

\begin{proposition}[Character-valued descent]
\label{prop:q4-sign-descent}
There is a nontrivial character
\[
 \chi:\Gamma_4\longrightarrow\{\pm1\}
\]
with
\begin{equation}
 \chi(J)=-1,
 \qquad \chi(JT)=-1,
 \qquad \chi(T)=+1,
 \label{eq:q4-sign-character-values}
\end{equation}
such that
\begin{equation}
 F(\gamma\tau)=\chi(\gamma)F(\tau),
 \quad
 W(\gamma\tau)=
 \begin{cases}W(\tau),&\chi(\gamma)=1,\\W(\tau)^{-1},&\chi(\gamma)=-1,
 \end{cases}
 \label{eq:q4-sign-action}
\end{equation}
and consequently
\begin{equation}
 \gamma^*g_4=g_4,
 \qquad
 a_4(\gamma\tau)=\chi(\gamma)a_4(\tau).
 \label{eq:q4-sign-descent}
\end{equation}
Thus the metric and zero network descend to the triangle orbifold, while $a_4$
descends as a section of the real line bundle associated with $\chi$.  It becomes
an ordinary scalar assignment on the index-two cover associated with
$\ker\chi$.
\end{proposition}

\begin{proof}
Both $F\circ\gamma$ and $F$ square to the same meromorphic function.  Their ratio
is therefore a constant in $\{\pm1\}$ on the connected surface $\HH$.  Composition
makes that sign a character.  The order-$2$ stabilizer $J$ acts by $z\mapsto-z$
in the local coordinate for which $F\sim z$, so $\chi(J)=-1$.  At the order-$4$
orbit, $JT$ acts by $z\mapsto\pm iz$ while $F\sim cz^{-2}$, and hence
$\chi(JT)=-1$.  Since $\chi(JT)=\chi(J)\chi(T)$, this gives $\chi(T)=+1$ and proves
\eqref{eq:q4-sign-character-values}.  Substitution in the Cayley transform gives
\eqref{eq:q4-sign-action}; logarithmic modulus then
sends $s$ to $\chi(\gamma)s$.  The metric is unchanged and the odd assignment
transforms by the same sign.
\end{proof}

This descent statement is one point at which the signed arithmetic history is
strictly richer than the unlabelled zero set.  It also fixes the global status of
the construction: on the orbifold the natural assignment is line-bundle valued,
not a silently single-valued real function.

\subsection{From zero histories to relative divisors}

The Hecke model supplies an arithmetic source for a zero network, but the phrase
``different paths to zero'' still has several inequivalent meanings.  We use the
following audit hierarchy.

\begin{center}
\small
\begin{tabular}{p{0.22\textwidth}p{0.67\textwidth}}
\toprule
Level & Equality being imposed \\
\midrule
raw word & literal equality of operation symbols and parentheses \\
marked history & equality of the sequential tree, accumulator, and operand slots \\
projective operator & equality after the Paper I evaluation map to
$\operatorname{PGL}_2(K)$ \\
endpoint & equality after applying an operator to a declared seed or parameter \\
prime over $K$ & equality inside one irreducible component defined over
the coefficient field $K$ \\
geometric sheet & equality only after extension to $\overline K$ or along a local
analytic branch \\
\bottomrule
\end{tabular}
\end{center}

The first four levels form successively coarser process data.  The last two are
not further syntactic quotients: they are decompositions of the endpoint equation.
Distinct operators can agree at one endpoint, while one $K$-irreducible component
can split into several geometric sheets permuted by Galois or monodromy.

Let $B$ be an integral parameter space over a characteristic-zero field $K$.  If
a marked history $h$ and a chosen seed produce a nonzero rational endpoint function
$\nu_h\in K(B)^\times$, define its codimension-one zero datum to be the zero part
of the principal divisor $(\nu_h)$.  The exceptional case $\nu_h\equiv0$ is not a
codimension-one divisor: the entire generic base is a zero condition and must be
recorded separately.  Operator-equivalent histories give the same endpoint
function, but the converse need not hold.  For a finite set $A$ of operator
classes, one possible encoding on a common regular domain is
\begin{equation}
 P_A(b,z)=\prod_{[h]\in A}\bigl(z-\nu_h(b)\bigr).
 \label{eq:finite-history-polynomial}
\end{equation}
The vanishing condition $P_A(b,0)=0$ records that at least one selected operator
history reaches zero.  Formula~\eqref{eq:finite-history-polynomial} is useful but not
canonical: it depends on the seed, the finite selection, multiplicities, and the
chosen operator quotient.

The following standard algebraic fact identifies the precise irreducible layer
once a monic arithmetic polynomial encoding a declared finite history selection
has been supplied.

\begin{theorem}[Arithmetic and geometric zero components]
\label{thm:relative-zero-divisor-decomposition}
Let $R$ be a normal finitely generated domain over a characteristic-zero field
$k$,
let $B=\Spec R$, let $K=\operatorname{Frac}(R)$, and let
\[
 P(z)=z^n+c_1z^{n-1}+\cdots+c_n\in R[z]
\]
have nonzero discriminant.  Write
\[
 \mathcal D_P=V(P)\subset\mathbb A^1_B.
\]
Then:
\begin{enumerate}
\item $\mathcal D_P$ is an effective relative Cartier divisor, finite flat of
degree $n$ over $B$;
\item the distinct irreducible factorization
\begin{equation}
 P=\prod_{i=1}^r P_i\quad\text{in }K[z]
 \label{eq:generic-history-factorization}
\end{equation}
is the unique decomposition of the generic zero divisor into $K$-prime
components; their closures are the horizontal prime components of the associated
Weil divisor on $\mathbb A^1_B$;
\item over a separable closure $K^{\mathrm{sep}}$, the roots of each $P_i$ form
one transitive $\Gal(K^{\mathrm{sep}}/K)$-orbit, although they may be several
geometric sheets;
\item on the open set
$B^\circ=B\setminus V(\operatorname{disc}P)$, the map
$\mathcal D_P^\circ\to B^\circ$ is finite \'etale.  If an embedding $k\hookrightarrow
\C$ is fixed, set $B_\C^\circ=B^\circ\times_k\C$ and
$\mathcal D_{P,\C}^\circ=\mathcal D_P^\circ\times_k\C$.  On each connected
component of $(B_\C^\circ)^{\mathrm{an}}$, the analytic covering components of
$(\mathcal D_{P,\C}^\circ)^{\mathrm{an}}$ are exactly the orbits of topological
monodromy on one fiber.
\item if $D_i$ denotes the horizontal prime closure belonging to $P_i$ and
$\widetilde D_i$ its normalization, then $\widetilde D_i\to B$ is finite and its
restriction over $B^\circ$ agrees with the corresponding finite \'etale covering
component.
\end{enumerate}
\end{theorem}

\begin{proof}
Because $P$ is monic, $R[z]/(P)$ is a free $R$-module with basis
$1,z,\ldots,z^{n-1}$; hence the zero scheme is finite flat of degree $n$.
The element $P$ is a non-zero-divisor in $R[z]$, so its vanishing defines an
effective Cartier divisor.  Unique factorization in the principal ideal domain
$K[z]$ gives \eqref{eq:generic-history-factorization}.  Taking closures of the
generic prime divisors gives all horizontal codimension-one components; monicity
precludes a vertical component.  Characteristic zero makes every $P_i$ separable,
and irreducibility is equivalent to transitivity of the absolute Galois action on
its roots.  After inverting the discriminant, $P$ and $P'$ generate the unit ideal
in the finite algebra, which is the finite-\'etale criterion.  The monodromy
assertion is the standard classification of covering components by their orbits.
For item~(5), schemes of finite type over a field are excellent, so the
normalization of each $D_i$ is finite.  Over $B^\circ$, the prime component is
already normal because it is finite \'etale over the normal base; hence
normalization changes nothing there
\cite{Hartshorne1977AlgebraicGeometry,GrothendieckRaynaud1971SGA1}.
\end{proof}

The theorem makes ``irreducible paths to zero'' field-relative.  A $K$-prime
component is an arithmetic unit of zero behavior; its geometric branches need not
be individually defined over $K$.  The discriminant records finite collisions and
ramification in this monic setting.  It does not detect the boundary, metric,
domain, or nonproper escape mechanisms included in the structural discriminant of
Section~\ref{sec:discriminants}.

\begin{remark}[Arithmetic reduction interface]
\label{rem:arithmetic-reduction-interface}
Suppose in addition that the family is given by an integral model over a number
ring.  Let $b$ be a closed point of the parameter base with finite residue field,
lying away from the discriminant and every declared bad-reduction locus.  Then the
degrees of the irreducible factors of the specialized polynomial $P_b$ are the
cycle lengths of geometric Frobenius on the corresponding finite \'etale fiber.
At a discriminant point, the local discriminant records failure of \'etaleness.
After normalizing in a splitting field, inertia records the ramified part, while
nonmaximality of an order or a bad integral presentation may contribute additional
discriminant factors.
Thus discriminant primes, inertia, and reduction factorization are genuinely
arithmetic zero data, not merely another drawing of complex monodromy.  Using
them for AEG histories still presupposes the functorial integral polynomial whose
construction is posed below \cite{Neukirch1999AlgebraicNumberTheory}.
\end{remark}

\subsection{What is proved, and what remains a program}

For the $q=4$ model, the arithmetic-to-geometric chain established in this paper
is
\begin{equation}
 \frac{\text{projective arithmetic histories over }K_4}
      {\text{operator equality}}
 \supset \Gamma_4
 \longrightarrow \beta^{-1}([0,1])
 \longrightarrow (g_4,a_4).
 \label{eq:q4-established-chain}
\end{equation}
If $e$ is one reference edge from an order-$2$ point to an order-$4$ point, the
edges are translates $\gamma e$.  Equality of two edge labels is equality of the
corresponding cosets of the edge stabilizer.  Four distinct incident cosets meet
at an order-$4$ endpoint.  Raw words representing the same $\gamma$ have already
been identified.  Thus the four prongs have an exact operator-coset meaning, but
not yet a canonical prime-divisor meaning.

\begin{remark}[Symmetry is not nonunique factorization]
The matrices of $\Gamma_4$ are defined over $\Z[\sqrt2]$, which is
norm-Euclidean and hence a unique factorization domain.  The arithmetic richness
proved by Theorem~\ref{thm:q4-hecke-zero-network} comes from noncommuting
projective operators, Hecke continued-fraction paths, and automorphic
uniformization; it does not come from nonunique factorization of elements.  If
``irreducible paths'' is intended specifically to mean inequivalent irreducible
factorizations, one must pass to a number ring or order with nontrivial
factorization theory, for example a nontrivial class group, and declare whether
elements, ideals, units, and order are being quotiented.  That extension is not a
claim of the present $q=4$ model.

For completeness, norm-Euclideanity follows directly.  Given
$x+y\sqrt2\in\Q(\sqrt2)$, choose $a,b\in\Z$ with
$|x-a|,|y-b|\le1/2$.  Then
\[
 \left|N\bigl((x-a)+(y-b)\sqrt2\bigr)\right|
 =\left|(x-a)^2-2(y-b)^2\right|\le\frac12<1.
\]
Applying this to a quotient gives Euclidean division by the absolute norm.
\end{remark}

\begin{problem}[History-to-divisor functor]
\label{prob:history-to-divisor-functor}
Construct a functorial rule
\[
 \mathfrak H\longmapsto
 (B_{\mathfrak H},P_{\mathfrak H},\mathcal D_{\mathfrak H})
\]
from a declared category of marked arithmetic histories to parameter spaces,
history polynomials, and relative zero divisors, satisfying all of the following:
\begin{enumerate}
\item marked-history equivalences and operator equalities act by specified
isomorphisms rather than by accidental duplicate factors;
\item ordinary admissibility, projective poles, and chart changes remain visible;
\item composition of histories has a compatible operation on divisors;
\item base change gives the expected Galois and analytic monodromy actions;
\item finite truncations admit a controlled locally finite limit when infinitely
many history classes are present.
\end{enumerate}
Determine whether the Hecke zero network is a natural infinite or automorphic
instance of this functor, rather than merely its operator-level precursor.
\end{problem}

This problem is the main arithmetic naturality gate of the paper.  The relative
divisor theorem does not solve it: that theorem starts after $P_{\mathfrak H}$ has
been given.  Conversely, the exact Hecke model does not solve it either: it
derives a zero network from a projective operator group, but it does not yet
associate a finite monic endpoint polynomial to arbitrary AEG histories.
